\part{SliceNet and the machine-learning pipeline}

\FigSliceNet

\section{Why machine learning?}

\begin{qa}{357}{What specific decision does SliceNet make?}
Given a flow's traffic shape (31 statistics over a window and the first 32 packets) at one vantage point, it predicts
its service class (healthcare, education, agriculture, general). In the loop, that prediction picks the HTB class a
flow from a multi-service device enters.
\end{qa}

\begin{qa}{358}{Why classify services if the traffic generators already know which service produced each flow?}
The generator's knowledge is \emph{ground truth for evaluation}, never given to the controller. In deployment the
network does not know the application. For single-institution UEs, the address is enough (host mode). SliceNet is for
devices whose address does not identify the service (shared kiosks, gateways, NATed hotspots).
\end{qa}

\begin{qa}{359}{Is SliceNet for the simulator, a future live deployment, or both?}
Both. It is evaluated offline on measured windows and runs live in emulation (in-loop, mixed UE, L3). A live real-network
deployment is future work (L4 absent).
\end{qa}

\begin{qa}{360}{Why not use source IPs, ports, application metadata or static rules?}
Those are better when available: the host and app modes \emph{are} such rules and are the defaults. Ports are
unreliable in practice (shared 443, NAT, dynamic ports) and DSCP is untrusted from UEs. ML is reserved for when
metadata is missing or untrustworthy.
\end{qa}

\begin{qa}{361}{When would a rule-based classifier beat ML?}
\begin{itemize}
\item When identity is in trusted metadata (the address of a clinic's device, an S-NSSAI in 5G);
\item when traffic is encrypted but endpoints are known;
\item when the deployment's traffic differs from training (ML degrades; rules do not);
\item when explainability is required.
\end{itemize}
In this thesis, host rules equal SliceNet in disaster (\NMixedDisasterHostSla{} vs \NMixedDisasterSlicenetSla) for the
mixed UE.
\end{qa}

\begin{qa}{362}{What does a classification error cost downstream?}
The flow gets another slice's allocation and AQM. Critical traffic misclassified as general loses its priority
share and floor; bulk traffic misclassified as healthcare can borrow healthcare capacity and inflate its queue.
\end{qa}

\begin{qa}{363}{Which errors are more dangerous: false critical, or missed critical?}
Missed critical traffic violates a life-relevant SLA. A false ``critical'' label lets bulk traffic consume the critical
slice, which can also violate it. Both are dangerous; the second can be capped by per-flow policing inside the critical
class.
\honest{Recall on healthcare is high (F1 \NSnRfHealthcare) but the asymmetric cost is not built into the decision
rule.}
\end{qa}

\begin{qa}{364}{What is the simplest baseline SliceNet should be compared with?}
The majority class (macro-F1 0.25 for 4 balanced classes); a direction-share rule (uplink-dominant $\rightarrow$
healthcare or agriculture); a port oracle (app mode, 100\% by construction).
\honest{Only the RSNSD-trained model and the RF/CNN pair were formally evaluated. Add a single-feature rule baseline.}
\end{qa}

\begin{qa}{365}{Why this model architecture?}
\why{Two representations: aggregate statistics (tabular) and the early packet sequence (temporal).}
\choice{An RF for the tabular statistics (strong with little data, robust to scale and collinearity) and a small 1D-CNN
fusing a packet-sequence branch with a shape MLP (local patterns such as frame bursts and request/response
alternation).}
\alt{Logistic regression (linear: too weak for multimodal classes), XGBoost (likely similar to RF), LSTM or
Transformer (more data-hungry).}
\trade{The RF won (\NSnRfFone{} vs \NSnCnnFone) at about 14$\times$ the inference latency of the CNN.}
\ev{Held-out-group evaluation, \NSnRfTest{} test windows.}{\Lthree}
\end{qa}

\begin{qa}{366}{Why not logistic regression, random forests, XGBoost, an MLP or another deep architecture?}
RF is used. Logistic regression is a good sanity baseline but cannot express the multimodal classes. XGBoost would
likely tie the RF. An MLP is subsumed by the CNN's shape branch. Larger deep models are not justified at
\NSnRows{} windows.
\honest{Logistic regression and XGBoost were not run.}
\end{qa}

\begin{qa}{367}{What evidence supports the architecture being suitable for the features?}
Empirically: the RF beats the CNN, which suggests the discriminative information is in aggregate statistics more than
in early packet order. The name ``SliceNet'' comes from prior work [1], and we do \emph{not} claim the CNN is best.
\end{qa}

\begin{qa}{368}{What is the input tensor or feature vector?}
RF: 31 shape features (\code{SHAPE_FEATURES}). CNN: a $3\times32$ tensor (signed size/1500, $\log(1+\text{IAT}_{ms})$,
TCP-control flag) plus the 31-feature vector.
\end{qa}

\begin{qa}{369}{What are the output classes?}
Healthcare, education, agriculture, general (closed set).
\end{qa}

\begin{qa}{370}{How are features encoded, normalised or scaled?}
RF: raw values, \texttt{NaN}$\rightarrow$0. CNN: $\text{sign}(x)\log(1+|x|)$ followed by a z-score with training
mean and std; sequences are pre-scaled by construction.
\end{qa}

\begin{qa}{371}{Where are the fitted preprocessing parameters stored?}
Inside the model artefacts: the RF pickle stores its column list; the CNN \code{.pt} stores the mean, std, column
order and state dict.
\end{qa}

\begin{qa}{372}{How is model versioning handled?}
By file names (\code{slicenet_baseline.pkl}, \code{slicenet_mininet.pkl}, \code{slicenet_cnn.pt}) plus metrics JSONs
with timestamps; models are gitignored and regenerated by \code{make models}.
\honest{There is no model registry and no hash link between a model and its training data.}
\end{qa}

\begin{qa}{373}{What happens if the model file is missing or incompatible?}
In-loop classification is disabled (\code{enabled=False}), so flows stay in host or default classes. The post-run
scorer falls back to the RSNSD model and records which one was used.
\honest{An incompatible pickle would raise at load.}
\end{qa}

\begin{qa}{374}{How do you prevent inference from using a different feature order than training?}
Both models re-index input columns to the stored order before predicting.
\end{qa}

\section{Training and model evaluation}

\begin{qa}{375}{What is the exact training procedure?}
\begin{enumerate}
\item Merge the collection runs.
\item Drop non-service labels.
\item Group split (25\% of run$\times$UE groups held out, seed 42).
\item RF: 200 trees, max depth 16, balanced class weights.
\item CNN: 80 epochs, Adam (lr $10^{-3}$, weight decay $10^{-4}$), batch 128, class-weighted cross-entropy.
\item Evaluate on the held-out groups.
\end{enumerate}
\code{classification/pipeline.py}.
\end{qa}

\begin{qa}{376}{Which loss function and optimiser, and why?}
Class-weighted cross-entropy (standard for multi-class, counteracts imbalance) with Adam (robust defaults for small
CNNs).
\end{qa}

\begin{qa}{377}{How are learning rate, batch size, epochs and regularisation selected?}
Fixed standard defaults, not tuned. Epochs were raised from 40 to 80 after the training loss had not plateaued.
Regularisation: batch norm, dropout 0.3, weight decay.
\honest{No hyper-parameter search, which favours neither model.}
\end{qa}

\begin{qa}{378}{How do you determine whether the model is overfitting?}
By the gap between training loss (0.23 at 80 epochs) and held-out performance, and by stability under group hold-out.
\honest{There was no validation curve. Add a validation split carved from the training groups.}
\end{qa}

\begin{qa}{379}{What is the role of the validation set?}
For tuning and early stopping. We had none because no tuning was done; the held-out groups are a pure test set.
\end{qa}

\begin{qa}{380}{When do you stop training?}
After a fixed 80 epochs. No early stopping.
\end{qa}

\begin{qa}{381}{How do you select a classification threshold, if any?}
Arg-max over four classes. There is no threshold and no reject option (Q396).
\end{qa}

\begin{qa}{382}{What are the accuracy, precision, recall, F1 and confusion-matrix results?}
\begin{itemize}
\item RF: macro-F1 \NSnRfFone, accuracy \NSnRfAcc. Per class: healthcare \NSnRfHealthcare, education
\NSnRfEducation, agriculture \NSnRfAgriculture, general \NSnRfGeneral.
\item CNN: macro-F1 \NSnCnnFone, accuracy \NSnCnnAcc.
\end{itemize}
Confusion matrices: \code{docs/figures/thesis/slicenet_confusion_*}. Most confusion is between education and general,
which share web request/response models.
\end{qa}

\begin{qa}{383}{Which metric matters most for your application, and why?}
Per-class recall of healthcare (a missed critical flow loses its SLA), then macro-F1 (every class counts equally,
regardless of window counts).
\end{qa}

\begin{qa}{384}{How does performance differ across classes?}
Healthcare and agriculture are near-perfect (uplink-dominant, distinctive periodic and video patterns). Education and
general are lower because both have downlink web and streaming traffic.
\end{qa}

\begin{qa}{385}{How does class imbalance affect the reported metrics?}
Accuracy would over-weight the frequent classes (low-rate IoT yields fewer windows). We report macro-F1 and use
balanced class weights.
\end{qa}

\begin{qa}{386}{What is the difference between macro-F1 and weighted-F1?}
Macro: the unweighted mean of per-class F1 (each class equal). Weighted: weighted by class support, so it is dominated
by large classes.
\end{qa}

\begin{qa}{387}{Why is accuracy alone insufficient?}
It hides class-specific failure. A classifier that never predicts the rare critical class can still score high
accuracy.
\end{qa}

\begin{qa}{388}{How do you quantify uncertainty in the reported results?}
\honest{Single group split, so there is no CI on F1 yet.}
Fix: repeated \code{GroupKFold} (5 folds $\times$ 3 seeds), reporting the mean $\pm$ CI. A bootstrap over test groups
gives a CI cheaply.
\end{qa}

\begin{qa}{389}{How much does performance vary across random seeds or data splits?}
Not measured for the final models (Q388). An indirect signal: the RF was 0.959 on the earlier 12-run dataset and
\NSnRfFone{} on the 18-run dataset; the CNN was 0.880 vs \NSnCnnFone. So the variation is about $\pm$0.01.
\end{qa}

\begin{qa}{390}{How do you decide whether an improvement is statistically meaningful?}
With repeated group CV, compare models with a paired test across folds (Wilcoxon or a corrected resampled t-test).
The RF-vs-CNN gap (about 0.08) is likely real given the $\pm$0.01 variation, but is untested.
\end{qa}

\begin{qa}{391}{What is the model's inference latency?}
CNN: \NInferCnnMs\,ms for a single flow, \NInferCnnBatchMs\,ms for a batch of 100 (CPU). RF: \NInferRfMs\,ms single,
dominated by thread-pool start-up (\code{n_jobs=-1}). Setting \code{n_jobs=1} would make the RF sub-millisecond.
\end{qa}

\begin{qa}{392}{What is the model's memory and CPU footprint?}
CNN: \NCnnParams{} parameters (\NCnnSizeKb\,KB). RF: \NRfSizeMb\,MB on disk. Both run on CPU only.
\end{qa}

\begin{qa}{393}{How does performance change when traffic patterns shift?}
The extreme shift we tested (RSNSD synthetic $\rightarrow$ measured) collapses performance to \NTrZero{}.
Within-family shift (new seeds, contexts and runs) keeps \NSnRfFone.
\honest{Moderate shifts (new apps, other codecs) are untested.}
\end{qa}

\begin{qa}{394}{How would you test generalisation to unseen scenarios, generators or network conditions?}
Leave-one-context-out, leave-one-application-out, train at 100\,Mb/s and test at 50\,Mb/s, train on Linux bridges and
test on OVS runs, and ultimately test on real traces.
\end{qa}

\begin{qa}{395}{How would you detect concept drift?}
Monitor the distribution of predicted classes and confidences per site, and feature drift (KS or PSI on key shape
features) against training. Periodic labelled audits. Trigger retraining.
\end{qa}

\begin{qa}{396}{What is the fallback when the model is uncertain or meets an unknown class?}
Today: none. Arg-max always assigns a class, but a flow below 5 packets stays in the default (general) class.
Proposed: if the maximum probability is below $\tau$, keep host or default treatment and re-try with more packets
(safe default).
\end{qa}

\section{Runtime inference and deployment}

\begin{qa}{397}{Is SliceNet actually invoked by the operational allocation and recovery path?}
Only with \code{classification=slicenet}, for unknown or mixed UEs: \code{OnlineSliceClassifier.step} runs every
control interval and installs per-flow u32 rules. Allocation and recovery never use it. They operate on slices, and
SliceNet only decides which slice a flow enters.
\end{qa}

\begin{qa}{398}{If not, what determines service identity on that path?}
By default, host rules (UE address $\rightarrow$ service), installed at QoS setup. In 3GPP mode, app rules (server and
port $\rightarrow$ slice type).
\end{qa}

\begin{qa}{399}{How would you integrate live feature extraction and inference into the control loop?}
It is already integrated for the mixed UE. Read the UE capture, group packets into flows, compute shape features and
sequences, predict, call \code{classify_flow} (\code{classification/online.py}). For scale: an incremental flow table
(not re-parsing the pcap), done at the gateway.
\end{qa}

\begin{qa}{400}{How is a predicted class mapped to a service policy?}
Class name $\rightarrow$ HTB classid (from \code{slices.yaml}) $\rightarrow$ a u32 filter on the 5-tuple, in both
directions, at the highest filter priority. The flow then gets that slice's allocation and AQM.
\end{qa}

\begin{qa}{401}{How do you prevent a wrong prediction from granting unlimited priority to non-critical traffic?}
The slice's HTB rate and ceiling bound it (no class is unlimited), and fq\_codel isolates flows within a class.
\honest{With borrowing up to link capacity, a misclassified bulk flow in the healthcare class could still crowd it.
Per-flow policing or a ceiling inside critical classes, plus a confidence threshold, are needed.}
\end{qa}

\begin{qa}{402}{Would you use confidence thresholds, hysteresis, multiple observations or policy constraints?}
Yes, all four:
\begin{itemize}
\item a confidence threshold;
\item k-of-n agreement across windows before installing;
\item stickiness (today the first decision sticks forever; it should be re-evaluated with hysteresis);
\item policy constraints (per-flow caps in critical classes).
\end{itemize}
\end{qa}

\begin{qa}{403}{How would you reconcile a prediction with trusted application metadata?}
Metadata wins (S-NSSAI in 5G, an authenticated app registration). ML fills gaps only, and disagreements are logged as
drift or abuse signals.
\end{qa}

\begin{qa}{404}{Where should inference run: edge gateway, central controller or a dedicated service?}
At the edge gateway (\code{s_rural} side), where the flows are seen and enforcement happens, which avoids exporting
traffic metadata. The model is small (\NCnnSizeKb\,KB CNN).
\end{qa}

\begin{qa}{405}{What is the latency budget for the classification-to-enforcement path?}
Decision after at least 5 packets, at the next control interval (up to 5\,s), plus milliseconds of inference and a
millisecond-scale \code{tc filter add}. For short critical flows this is too slow; metadata or host rules must cover
them.
\end{qa}

\begin{qa}{406}{What if inference is slower than the monitoring interval?}
It runs synchronously in the loop, so a slow model would delay the tick. The wall-clock-aligned scheduler would not
accumulate drift, but telemetry would be late. Mitigation: run inference in a worker thread with the latest-result
pattern, like the probes.
\end{qa}

\begin{qa}{407}{How would you deploy a new model without interrupting service?}
Load it next to the old one, run it in shadow mode (log predictions, do not act), compare agreement and confidence,
then atomically swap the reference. Existing flow rules persist.
\end{qa}

\begin{qa}{408}{How would you roll back a model that degrades performance?}
Keep the previous artefact. Watch per-slice SLA and confidence metrics; swap back on regression. Flow rules installed
by the bad model can be flushed.
\end{qa}

\begin{qa}{409}{What monitoring is needed to detect inference failures or distribution shift?}
Inference errors and latency, the prediction class mix vs expected, confidence histograms, feature drift statistics,
and the downstream SLA of classified flows.
\end{qa}

\begin{qa}{410}{How would you evaluate the classifier on real network traffic?}
Labelled ground truth from consenting institutions (known devices or apps), header-only capture, offline evaluation
first, then shadow mode in the loop before acting.
\end{qa}

\FigTransfer
